\begin{abstract}
Agent Skills combine persistent natural-language instructions with code,
configuration, and resources, allowing an untrusted package to shape an
agent's reasoning while inheriting its access to files, credentials, networks,
and processes. We study whether malicious Skills can be filtered before any
target content is installed, imported, or executed. \emph{Read Before You Run}
extracts bounded high-level descriptions, instruction semantics, normalized
sensitive operations, and cross-file Python call and taint paths; a
schema-constrained reviewer returns \textsc{Pass}, \textsc{Review}, or
\textsc{Block} with source-grounded evidence and explicit unknowns. On a
1,000-package source-disjoint benchmark, however, the complete method reaches
only 57/491 (11.61\%) malicious recall and 20.77\% F1 on 977 common successful
samples. Its 1/486 (0.21\%) false-positive rate is low, but a one-call document
baseline attains 36.15\% F1 and significantly more correct predictions
($p=3.17\times10^{-12}$). Removing the independently extracted high-level
function also does not hurt detection. The static front end nevertheless scans
1,000 current community Skills without failure, producing auditable candidates
without executing them. These results delimit, rather than overstate, the
current capability of zero-execution Skill screening: structured evidence
supports low-FPR triage and inspection, but does not yet improve source-disjoint
detection over direct semantic review.
\end{abstract}
